\section{Of which shape? Orthomorphisms and cycle types}\label{sec:cycle-types}

A complete mapping is a permutation of $G$, so it has a cycle type. It is more convenient to work with a closely related notion.

\begin{definition}\label{def:orthomorphism}
A bijection $\psi\colon G\to G$ is an \term{orthomorphism} if $x\mapsto x^{-1}\psi(x)$ is also a bijection. Equivalently, choosing in row $x$ of the Cayley table the cell with symbol $\psi(x)$ gives a transversal~\cite{wan07}.
\end{definition}

For example, the spade transversal of Ozanam's square, read as a permutation of $\Z_2\times\Z_2$, fixes $00$ and cycles $01\to10\to11\to01$ (Figure~\ref{fig:ortho-cycle}).

\begin{figure}[ht]
\centering
\begin{tikzpicture}[>=Latex, every node/.style={font=\normalsize}]
    \node[overtex, minimum size=8mm, inner sep=0pt] (e) at (0,0) {$00$};
    \draw[->, thick, BlkConjecture] (e) edge[loop above, looseness=6] (e);
    \node[overtex, minimum size=8mm, inner sep=0pt] (a) at (2.2,0.9) {$01$};
    \node[overtex, minimum size=8mm, inner sep=0pt] (b) at (3.5,-0.9) {$10$};
    \node[overtex, minimum size=8mm, inner sep=0pt] (c) at (0.9,-0.9) {$11$};
    \draw[->, thick, BlkConjecture] (a) to[bend left=20] (b);
    \draw[->, thick, BlkConjecture] (b) to[bend left=20] (c);
    \draw[->, thick, BlkConjecture] (c) to[bend left=20] (a);
  \end{tikzpicture}

\caption{The spade transversal as a permutation of $\Z_2\times\Z_2$.}
\label{fig:ortho-cycle}
\end{figure}

Prescribed cycle types have long-standing applications, as recorded by M\"uyesser~\cite{muy23}. A single cycle through all non-identity elements appears in Ringel's solution of the Heawood map-colouring problem, and orthomorphisms whose non-identity elements fall into disjoint $6$-cycles give cyclic Steiner triple systems.

\begin{conjecture}[Friedlander--Gordon--Tannenbaum~\cite{fgt81}]\label{conj:fgt}
Let $G$ be an abelian group of order $n$ satisfying the Hall--Paige condition, and let $k\geq2$ with $k\mid n-1$. Then some orthomorphism of $G$ fixes $0$ and permutes the other elements in disjoint $k$-cycles.
\end{conjecture}

Figure~\ref{fig:ortho-cycle} is the case $G=\Z_2\times\Z_2$, $k=3$. Friedlander, Gordon and Tannenbaum~\cite{fgt81} proved the conjecture for $|G|\leq15$ and for abelian $p$-groups with $p\geq3$. Even the case $k=3$ with $G$ cyclic remained open for forty years.

\begin{theorem}[M\"uyesser~\cite{muy23}]\label{thm:fgt}
The Friedlander--Gordon--Tannenbaum conjecture holds for all sufficiently large abelian groups.
\end{theorem}

\topic{Cycle types as rainbow cycle factors}

Let $\operatorname{Cay}(G,G\setminus\{e\})$ denote the complete directed Cayley graph in which the arc $x\to xs$ has colour $s$; equivalently, the arc $x\to y$ has colour $x^{-1}y$.

\begin{proposition}\label{prop:cyc}
Orthomorphisms $\psi$ of $G$ with $\psi(e)=e$ that permute $G\setminus\{e\}$ in cycles of lengths $k_1,k_2,\dots$ correspond exactly to rainbow spanning families of vertex-disjoint directed cycles of lengths $k_1,k_2,\dots$ on $G\setminus\{e\}$ in $\operatorname{Cay}(G,G\setminus\{e\})$.
\end{proposition}

\begin{proof}
Given $\psi$, use the arcs $x\to\psi(x)$ for $x\neq e$. Their colours $x^{-1}\psi(x)$ are distinct because $\psi$ is an orthomorphism, and none of them is $e$, because $e^{-1}\psi(e)=e$ is already taken. Conversely, a rainbow cycle factor on $G\setminus\{e\}$ uses $|G|-1$ distinct non-identity colours, that is, each non-identity colour exactly once. Setting $\psi(x)$ to be the successor of $x$ on its cycle and $\psi(e)=e$, the map $x\mapsto x^{-1}\psi(x)$ takes every value exactly once.
\end{proof}

A single cycle of length $|G|-1$ is one of Ringel's R-sequencings, and a factor of triangles is the case $k=3$ of the Friedlander--Gordon--Tannenbaum conjecture.
